\subsubsection{Conditioning of the Estimate}\label{subsubsec:condicionamento-R}


The singular values ​​of a matrix $\mathbf X$ are the square roots of the eigenvalues ​​of $\mathbf X\mathbf X^H$ and measure how much the matrix expands or reduces a vector in each direction. The rank of a matrix is ​​the number of its non-zero singular values, and measures the amount of non-redundant information in matrices. A sum of $N$ matrices of rank 1 has a rank of at most $N$ \cite{tse-viswanath}. Each term $\mathbf{n}^i(\mathbf{n}^i)^H$ of \eqref{eq:Rn_est} has rank 1, because all its columns are multiples of $\mathbf{n}^i$. Therefore:
% FONTE: Tse e Viswanath, Seção 7.1.1, pp. 291-292, Eqs. 7.2 a 7.4.

\[
\operatorname{rank}(\widehat{\mathbf R}_n) \leq \min(M,N)
\]

A matrix is ​​considered singular when its rank is strictly less than its dimension. Singular matrices do not have an inverse. Thus, $\widehat{\mathbf R}_n$ has no inverse when $N<M$, and $\hat{\mathbf R}_{cc}$ has no inverse when $N<M_c$.


The number of antennas in a massive MIMO system can exceed the number of samples in a coherence block (Table~\ref{tab:n-maximo}). Ledoit-Wolf shrinkage renders the estimate invertible. It blends $\widehat{\mathbf R}_n$ with a multiple of the identity matrix, using a weight calculated from the samples themselves \cite{ledoit-wolf}.
% FONTE: N da Tabela tab:n-maximo (24 a 2400 REs de DM-RS por PRB, 0 nos HST).
% FONTE: Ledoit e Wolf (2004). Simulador: m5_noise.shrink.

The condition number of a matrix is ​​the ratio between its largest and smallest singular values ​​\cite{tse-viswanath}. The error in matrix inversion using a low number of bits increases with the condition number, and this number is linked to the ratio of interference to noise \cite{osinsky2023}. In this work, this ratio is measured via the IoTh (Section~\ref{subsec:ioth}).
% FONTE: Tse e Viswanath, Seção 7.1.2, p. 295; Osinsky et al. 2023, p. 1603.
% FONTE: Tse e Viswanath, Seção 3.3.2, p. 75, nota 9, Eq. 3.83; Eq. 7.2, p. 292.
